% questions.tex —— 2014 年·全国大纲卷·文科（广西） · 题目内容片段（无答案）
% 由 tools/generate_projects.py 依来源总账自 2014/2014_数学真题_大纲文_试题.md 生成，勿手改（改源 MD 后重跑）


\begin{gknotice}
\textbf{来源说明}：本题面依据固定源 PDF 的两页 A3 题面逐题转录，原卷印刷页码为 871—872。第 871 面为选择题 1—12、填空题 13—16、解答题 17—19；第 872 面为解答题 20—22。

\textbf{图形说明}：第 19 题三棱柱图由项目内 TikZ 矢量源重绘；其余各题原页无图。

\textbf{解析归属}：固定源不附答案或解析；配套答案与解析均由本库按本题面独立复算。
\end{gknotice}

\gkpart{一、选择题}

\begin{question}
\gkti{1}设集合 $M=\{1,2,4,6,8\}$，$N=\{1,2,3,5,6,7\}$，则 $M\cap N$ 中元素的个数为\mbox{（　　）}

（A）$2$；（B）$3$；（C）$5$；（D）$7$。

\end{question}
\begin{question}
\gkti{2}已知角 $\alpha$ 的终边经过点 $(-4,3)$，则 $\cos\alpha=$\mbox{（　　）}

（A）$\dfrac45$；（B）$\dfrac35$；（C）$-\dfrac35$；（D）$-\dfrac45$。

\end{question}
\begin{question}
\gkti{3}不等式组 $\begin{cases}x(x+2)>0,\\|x|<1\end{cases}$ 的解集为\mbox{（　　）}

（A）$\{x\mid -2<x<-1\}$；（B）$\{x\mid -1<x<0\}$；（C）$\{x\mid 0<x<1\}$；（D）$\{x\mid x>1\}$。

\end{question}
\begin{question}
\gkti{4}已知正四面体 $ABCD$ 中，$E$ 是 $AB$ 的中点，则异面直线 $CE$ 与 $BD$ 所成角的余弦值为\mbox{（　　）}

（A）$\dfrac16$；（B）$\dfrac{\sqrt3}{6}$；（C）$\dfrac13$；（D）$\dfrac{\sqrt3}{3}$。

\end{question}
\begin{question}
\gkti{5}函数 $y=\ln(\sqrt[3]{x}+1)$（$x>-1$）的反函数是\mbox{（　　）}

（A）$y=(1-\mathrm e^x)^3$（$x>-1$）；（B）$y=(\mathrm e^x-1)^3$（$x>-1$）；（C）$y=(1-\mathrm e^x)^3$（$x\in\mathbf R$）；（D）$y=(\mathrm e^x-1)^3$（$x\in\mathbf R$）。

\end{question}
\begin{question}
\gkti{6}已知 $\boldsymbol a$、$\boldsymbol b$ 为单位向量，其夹角为 $60^\circ$，则 $(2\boldsymbol a-\boldsymbol b)\cdot\boldsymbol b=$\mbox{（　　）}

（A）$-1$；（B）$0$；（C）$1$；（D）$2$。

\end{question}
\begin{question}
\gkti{7}有 6 名男医生、5 名女医生，从中选出 2 名男医生、1 名女医生组成一个医疗小组，则不同的选法共有\mbox{（　　）}

（A）$60$ 种；（B）$70$ 种；（C）$75$ 种；（D）$150$ 种。

\end{question}
\begin{question}
\gkti{8}设等比数列 $\{a_n\}$ 的前 $n$ 项和为 $S_n$，若 $S_2=3$，$S_4=15$，则 $S_6=$\mbox{（　　）}

（A）$31$；（B）$32$；（C）$63$；（D）$64$。

\end{question}
\begin{question}
\gkti{9}已知椭圆 $C:\dfrac{x^2}{a^2}+\dfrac{y^2}{b^2}=1$（$a>b>0$）的左、右焦点分别为 $F_1$、$F_2$，离心率为 $\dfrac{\sqrt3}{3}$，过 $F_2$ 的直线 $l$ 交 $C$ 于 $A$、$B$ 两点，若 $\triangle AF_1B$ 的周长为 $4\sqrt3$，则 $C$ 的方程为\mbox{（　　）}

（A）$\dfrac{x^2}{3}+\dfrac{y^2}{2}=1$；（B）$\dfrac{x^2}{3}+y^2=1$；（C）$\dfrac{x^2}{12}+\dfrac{y^2}{8}=1$；（D）$\dfrac{x^2}{12}+\dfrac{y^2}{4}=1$。

\end{question}
\begin{question}
\gkti{10}正四棱锥的顶点都在同一球面上，若该棱锥的高为 $4$，底面边长为 $2$，则该球的表面积为\mbox{（　　）}

（A）$\dfrac{81\pi}{4}$；（B）$16\pi$；（C）$9\pi$；（D）$\dfrac{27\pi}{4}$。

\end{question}
\begin{question}
\gkti{11}双曲线 $C:\dfrac{x^2}{a^2}-\dfrac{y^2}{b^2}=1$（$a>0$，$b>0$）的离心率为 $2$，焦点到渐近线的距离为 $\sqrt3$，则 $C$ 的焦距等于\mbox{（　　）}

（A）$2$；（B）$2\sqrt2$；（C）$4$；（D）$4\sqrt2$。

\end{question}
\begin{question}
\gkti{12}奇函数 $f(x)$ 的定义域为 $\mathbf R$，若 $f(x+2)$ 为偶函数，$f(1)=1$，则 $f(8)+f(9)=$\mbox{（　　）}

（A）$-2$；（B）$-1$；（C）$0$；（D）$1$。

\end{question}
\gkpart{二、填空题}

\begin{question}
\gkti{13}$(x-2)^6$ 的展开式中 $x^3$ 的系数为\gkblank{}。（用数字作答）

\end{question}
\begin{question}
\gkti{14}函数 $y=\cos2x+2\sin x$ 的最大值为\gkblank{}。

\end{question}
\begin{question}
\gkti{15}设 $x$、$y$ 满足约束条件
\nopagebreak[4]
\[\begin{cases}
x-y\geqslant 0,\\
x+2y\leqslant 3,\\
x-2y\leqslant 1,
\end{cases}\]
则 $z=x+4y$ 的最大值为\gkblank{}。

\end{question}
\begin{question}
\gkti{16}直线 $l_1$ 和 $l_2$ 是圆 $x^2+y^2=2$ 的两条切线，若 $l_1$ 与 $l_2$ 的交点为 $(1,3)$，则 $l_1$ 与 $l_2$ 的夹角的正切值等于\gkblank{}。

\end{question}
\gkpart{三、解答题}

\begin{question}
\gkti{17}数列 $\{a_n\}$ 满足 $a_1=1$，$a_2=2$，$a_{n+2}=2a_{n+1}-a_n+2$。

（1） 设 $b_n=a_{n+1}-a_n$，证明 $\{b_n\}$ 是等差数列；

（2） 求数列 $\{a_n\}$ 的通项公式。

\end{question}
\begin{question}
\gkti{18}$\triangle ABC$ 的内角 $A$、$B$、$C$ 的对边分别为 $a$、$b$、$c$，已知 $3a\cos C=2c\cos A$，$\tan A=\dfrac13$，求 $B$。

\Needspace{0.4\textheight}
\end{question}
\begin{question}
\gkti{19}如图，三棱柱 $ABC-A_1B_1C_1$ 中，点 $A_1$ 在平面 $ABC$ 内的射影 $D$ 在 $AC$ 上，$\angle ACB=90^\circ$，$BC=1$，$AC=CC_1=2$。
\nopagebreak[4]
\begin{center}% ============================================================
% f2014_outline_arts_q19_prism.tex
% 2014 全国大纲卷（文科）第 19 题：三棱柱 ABC-A1B1C1
% A1 在平面 ABC 内的射影 D 在 AC 上，∠ACB=90°。
% 来源 PDF：普通高考/2014/2014大纲文(广西).pdf
% 原页：第 1 页（A3 横排）第 3 栏，印面 871，第 19 题题干下方
% ============================================================
\begin{tikzpicture}[>=Stealth, line width=0.7pt, font=\small,
  line cap=round, line join=round]
  \coordinate (A) at (0,0);
  \coordinate (C) at (2.55,0.85);
  \coordinate (B) at (4.35,0.65);
  \coordinate (D) at (0.95,0.3167);
  \coordinate (A1) at (0.72,3.15);
  \coordinate (C1) at (3.27,4.00);
  \coordinate (B1) at (5.07,3.80);

  % 可见棱
  \draw (A) -- (B) -- (B1) -- (C1) -- (A1) -- cycle;
  \draw (A) -- (A1);
  \draw (A1) -- (B1);
  \draw (A1) -- (B);

  % 不可见棱与辅助线
  \draw[dashed] (A) -- (C) -- (B);
  \draw[dashed] (C) -- (C1);
  \draw[dashed] (A1) -- (D);
  \draw[dashed] (A1) -- (C);

  \fill (A) circle (1.1pt) node[below left] {$A$};
  \fill (B) circle (1.1pt) node[right] {$B$};
  \fill (C) circle (1.1pt) node[above left] {$C$};
  \fill (D) circle (1.1pt) node[below right] {$D$};
  \fill (A1) circle (1.1pt) node[left] {$A_1$};
  \fill (C1) circle (1.1pt) node[above] {$C_1$};
  \fill (B1) circle (1.1pt) node[above right] {$B_1$};
\end{tikzpicture}
\end{center}
\nopagebreak[4]
（1） 证明：$AC_1\perp A_1B$；

（2） 设直线 $AA_1$ 与平面 $BCC_1B_1$ 的距离为 $\sqrt3$，求二面角 $A_1-AB-C$ 的大小。

\end{question}
\begin{question}
\gkti{20}设每个工作日甲、乙、丙、丁 4 人需使用某种设备的概率分别为 $0.6$、$0.5$、$0.5$、$0.4$，各人是否需使用设备相互独立。

（1） 求同一工作日至少 3 人需使用设备的概率；

（2） 实验室计划购买 $k$ 台设备供甲、乙、丙、丁使用，若要求“同一工作日需使用设备的人数大于 $k$”的概率小于 $0.1$，求 $k$ 的最小值。

\end{question}
\begin{question}
\gkti{21}函数 $f(x)=ax^3+3x^2+3x$（$a\neq0$）。

（1） 讨论函数 $f(x)$ 的单调性；

（2） 若函数 $f(x)$ 在区间 $(1,2)$ 是增函数，求 $a$ 的取值范围。

\end{question}
\begin{question}
\gkti{22}已知抛物线 $C:y^2=2px$（$p>0$）的焦点为 $F$，直线 $y=4$ 与 $y$ 轴的交点为 $P$，与 $C$ 的交点为 $Q$，且 $|QF|=\dfrac54|PQ|$。

（1） 求 $C$ 的方程；

（2） 过 $F$ 的直线 $l$ 与 $C$ 相交于 $A$、$B$ 两点，若 $AB$ 的垂直平分线 $l'$ 与 $C$ 相交于 $M$、$N$ 两点，且 $A$、$M$、$B$、$N$ 四点在同一圆上，求 $l$ 的方程。

\end{question}